\label{sec:proposed}

\subsection{Central Hypothesis}

We hypothesise that session-level feature vectors derived exclusively from Windows Event Log
event identifiers, timestamps, and metadata fields contain sufficient discriminative
information to classify smart grid host sessions as benign or malicious with macro
F1-score $\geq 0.95$ under realistic class imbalance ($5:1$ benign-to-malicious ratio).

\subsection{Research Objectives}

\begin{enumerate}
  \item \textbf{Synthetic Data Fidelity} (RO1): Design a generator that reproduces the
        statistical distributions of event-ID sequences observed in published ICS incident
        reports, covering at least four distinct attack categories. Success criterion:
        generated sessions are indistinguishable from real logs by a domain expert
        assessment on a 50-session blind review.

  \item \textbf{Feature Informativeness} (RO2): Engineer session-level features that
        individually achieve mutual information $> 0.01$ with the session label on the
        training set. Success criterion: no feature group removal in the ablation study
        reduces macro F1 by less than 2 percentage points.

  \item \textbf{Detection Performance} (RO3): Achieve macro F1-score $\geq 0.95$ and
        ROC-AUC $\geq 0.98$ on the held-out test partition under 5-fold stratified
        cross-validation. Success criterion: metrics reported with 95\% bootstrap confidence
        intervals.

  \item \textbf{Operational Applicability} (RO4): Ensure inference latency per session is
        $<10$~ms on commodity hardware (single CPU core), making real-time deployment
        feasible alongside existing SIEM pipelines. Success criterion: latency measured over
        1{,}000 sessions, median and 99th percentile reported.
\end{enumerate}

\subsection{Novelty and Scope}

The proposed work is novel in four respects. First, it targets Windows Event Logs at the
\emph{host level} within the OT/IT convergence zone—a layer confirmed absent from existing
ICS IDS literature \cite{mitchell2014survey,conti2021survey}. Second, the session
abstraction bridges the gap between raw event streams and sequence-level attack pattern
recognition, capturing temporal burst structure through entropy and time-delta statistics
without the inference-latency cost of deep sequence models \cite{du2017deeplog}. Third, the
open-source synthetic generator—including adversarially-timed evasive attack variants—enables
community replication without access to sensitive operational data, addressing the data
availability barrier identified by \cite{nafees2023}. Fourth, the systematic ablation study
provides deployment guidance: practitioners with constrained feature coverage (e.g.,
no timestamp resolution) can prioritise event-count features based on quantified delta-F1.
The scope is bounded to binary (benign/malicious) session classification; multi-class
attack-type attribution per MITRE ATT\&CK for ICS \cite{mitre2020attack} is left for
future work.
